\documentclass[11pt]{article}
\usepackage[hmargin=2.5cm,vmargin=2.5cm]{geometry}
\usepackage{amsmath,amssymb}
\usepackage[colorlinks=true,linkcolor=blue,urlcolor=blue]{hyperref}

\title{The Household's Problem:\\ FOCs, the Human-Capital Margin, and the SDF}
\author{Spending-efficiency model --- internal note}
\date{\today}

\begin{document}
\maketitle

\noindent
This note derives the household block of the model
(\texttt{models/model\_block.modpart}, households section): the six first-order
conditions, the work-vs-schooling arbitrage, and the stochastic discount factor.
Notation is the paper's throughout. The derivation is in levels; the aside at the
end records which conditions pick up growth factors when the model is
stationarized.

\section{The Household's Problem}

There is a representative household that consumes, supplies labor, spends time
building human capital, and accumulates private capital and government bonds. Its
problem is:
\begin{equation*}
    \max_{C_t,\,L_t,\,E_t,\,K_t,\,B_t,\,H_t} \;
    \mathbb{E}_0 \sum_{t=0}^{\infty} \beta^{t}
    \left[\ln C_t - \omega\,\frac{\left(L_t + E_t\right)^{1+\varphi}}{1+\varphi}\right]
\end{equation*}
subject to the flow budget constraint,
\begin{equation}\label{eq:budget}
    \left(1+\tau_t^{c}\right) C_t + I_t + B_t
    = \left(1-\tau_t^{w}\right) w_t\, N_t + r_t^{k} K_{t-1}
    + \frac{R_{t-1}}{\Pi_t}\, B_{t-1} + T_t + \Omega_t,
\end{equation}
the private-capital law of motion $K_t = (1-\delta)K_{t-1} + I_t$, the
human-capital law of motion,
\begin{equation}\label{eq:hlom}
    H_t = \left(1-\delta^{h}\right) H_{t-1}
    + \chi\, E_t^{\gamma} \left(\frac{K_{t-1}^{GE}}{\Lambda_{t-1}}\right)^{\mu},
\end{equation}
and the definition of effective labor, $N_t = H_{t-1} L_t$. The household takes
prices, taxes, transfers, and the rebated profits $\Omega_t$ as given. Note two
things about the time margins. First, labor $L_t$ and education time $E_t$ enter
the disutility term jointly, so they compete for the same time endowment. Second,
the human-capital stock enters effective labor with a lag: an hour of schooling
today raises earnings capacity only from tomorrow on.

Substitute investment $I_t = K_t - (1-\delta)K_{t-1}$ and effective labor
$N_t = H_{t-1}L_t$ out of \eqref{eq:budget}, and attach the multiplier $\lambda_t$
to the budget constraint and $\lambda_t^{H}$ to \eqref{eq:hlom}. A Lagrangian for
the household is:
\begin{equation*}
\begin{aligned}
    \mathbb{L} = \mathbb{E}_0 \sum_{t=0}^{\infty} \beta^{t} \Big\{\,
    & \ln C_t - \omega\,\frac{\left(L_t + E_t\right)^{1+\varphi}}{1+\varphi} \\
    & + \lambda_t \Big[\left(1-\tau_t^{w}\right) w_t H_{t-1} L_t
      + r_t^{k} K_{t-1} + \tfrac{R_{t-1}}{\Pi_t} B_{t-1} + T_t + \Omega_t \\
    & \qquad\;\; - \left(1+\tau_t^{c}\right) C_t
      - \big(K_t - (1-\delta) K_{t-1}\big) - B_t \Big] \\
    & + \lambda_t^{H} \Big[\left(1-\delta^{h}\right) H_{t-1}
      + \chi\, E_t^{\gamma} \big(K_{t-1}^{GE}/\Lambda_{t-1}\big)^{\mu}
      - H_t \Big] \Big\}.
\end{aligned}
\end{equation*}

\section{First-Order Conditions}

The FOC are:
\begin{align}
    \frac{\partial\mathbb{L}}{\partial C_t} = 0
    &\;\Leftrightarrow\;
    \frac{1}{C_t} = \lambda_t \left(1+\tau_t^{c}\right), \label{eq:focC}\\
    \frac{\partial\mathbb{L}}{\partial B_t} = 0
    &\;\Leftrightarrow\;
    \lambda_t = \beta\,\mathbb{E}_t\!\left[\lambda_{t+1}\,
    \frac{R_t}{\Pi_{t+1}}\right], \label{eq:focB}\\
    \frac{\partial\mathbb{L}}{\partial K_t} = 0
    &\;\Leftrightarrow\;
    \lambda_t = \beta\,\mathbb{E}_t\!\left[\lambda_{t+1}
    \left(r_{t+1}^{k} + 1 - \delta\right)\right], \label{eq:focK}\\
    \frac{\partial\mathbb{L}}{\partial L_t} = 0
    &\;\Leftrightarrow\;
    \omega\left(L_t + E_t\right)^{\varphi}
    = \lambda_t \left(1-\tau_t^{w}\right) w_t H_{t-1}, \label{eq:focL}\\
    \frac{\partial\mathbb{L}}{\partial E_t} = 0
    &\;\Leftrightarrow\;
    \omega\left(L_t + E_t\right)^{\varphi}
    = \lambda_t^{H}\, \chi\gamma\, E_t^{\gamma-1}
    \big(K_{t-1}^{GE}/\Lambda_{t-1}\big)^{\mu}, \label{eq:focE}\\
    \frac{\partial\mathbb{L}}{\partial H_t} = 0
    &\;\Leftrightarrow\;
    \lambda_t^{H} = \beta\,\mathbb{E}_t\!\left[\lambda_{t+1}
    \left(1-\tau_{t+1}^{w}\right) w_{t+1} L_{t+1}
    + \lambda_{t+1}^{H}\left(1-\delta^{h}\right)\right]. \label{eq:focH}
\end{align}
For \eqref{eq:focK}, note that $K_t$ appears twice in the Lagrangian: bought today
(through the substituted investment, at utility price $\lambda_t$) and carried
into tomorrow's income and resale (at $\lambda_{t+1}(r^k_{t+1}+1-\delta)$). For
\eqref{eq:focH} the structure is the same: an extra unit of $H_t$ raises
tomorrow's effective labor --- and hence after-tax labor income
$\lambda_{t+1}(1-\tau^w_{t+1})w_{t+1}L_{t+1}$ --- and survives into the future
stock at rate $1-\delta^{h}$, valued at $\lambda^{H}_{t+1}$.

\eqref{eq:focC} identifies $\lambda_t$ as the marginal utility of wealth.
\eqref{eq:focB}--\eqref{eq:focK} are the bond and capital Euler equations
(\texttt{lambda = betta*(lambda(+1)/g(+1)*R/PI(+1))} and
\texttt{1 = betta*(lambda(+1)/lambda/g(+1)*(1-delta+rk(+1)))} in the model; the
$g_{t+1}$ is the stationarization --- see the aside). \eqref{eq:focL} is the
labor-supply condition (\texttt{omega*(L+E)\^{}varphi = lambda*w*H(-1)*(1-tauw)}):
an hour of work supplies $H_{t-1}$ efficiency units paid $w_t$ each, so its
utility value is the after-tax hourly wage times the marginal utility of wealth.
\eqref{eq:focE} and \eqref{eq:focH} are the model equations
\texttt{omega*(L+E)\^{}varphi = lambda\_H*chiH*gamma*E\^{}(gamma-1)*
(Kge(-1))\^{}(mu)} and \texttt{lambda\_H = betta*(lambda(+1)*(1-tauw(+1))*w(+1)*
L(+1)+lambda\_H(+1)*(1-deltaH))}.

\section{The Work-vs-Schooling Arbitrage}

Because $L_t$ and $E_t$ draw on the same time endowment, their marginal
disutilities coincide --- the left-hand sides of \eqref{eq:focL} and
\eqref{eq:focE} are identical. Equating the right-hand sides:
\begin{equation}\label{eq:arbitrage}
    \underbrace{\lambda_t \left(1-\tau_t^{w}\right) w_t
    H_{t-1}}_{\text{return to an hour of work}}
    \;=\;
    \underbrace{\lambda_t^{H}\, \chi\gamma\, E_t^{\gamma-1}
    \big(K_{t-1}^{GE}/\Lambda_{t-1}\big)^{\mu}}_{\text{return to an hour of
    schooling}}.
\end{equation}
\eqref{eq:arbitrage} says the household allocates time until the after-tax wage
of the marginal hour equals the shadow value of the human capital that hour of
schooling builds. This is the margin the public human-capital experiments work
through: a higher $K^{GE}$ raises the marginal product of schooling time, tilting
\eqref{eq:arbitrage} toward education, and the resulting $H$ raises effective
labor with a lag.

\section{The Stochastic Discount Factor}

The household's valuation of a real payoff at $t+1$ per unit at $t$ is the ratio
of marginal utilities of \emph{income}. From \eqref{eq:focC}, a unit of income at
$t$ buys $1/(1+\tau^c_t)$ units of consumption, so:
\begin{equation}\label{eq:sdf}
    \mathit{SDF}_{t+1}
    = \beta\,\frac{\lambda_{t+1}\left(1+\tau_{t+1}^{c}\right)}
    {\lambda_t\left(1+\tau_t^{c}\right)}
    = \beta\,\frac{C_t}{C_{t+1}},
\end{equation}
where the second equality substitutes \eqref{eq:focC}; it is exact, with no
constant-tax assumption needed. \eqref{eq:sdf} is the model equation
\texttt{SDF = betta*lambda*(1+tauc)/(lambda(-1)*(1+tauc(-1)))}. The SDF is
defined by household preferences but used elsewhere: by the technology adopters
(\texttt{creationAdoption.tex}) and, in the $\beta^{s}\lambda_{t+s}$ form, by the
price-setting retailers (\texttt{pricingEquations.tex}).

\section{Aside: Which Conditions Pick Up $g$ When Stationarized}

Detrending maps $\lambda_t \mapsto \hat\lambda_t/\Lambda_t$ (marginal utility
falls as the economy grows) and leaves $H_t$, $L_t$, and $E_t$ untouched (all
stationary). The consequences, condition by condition:
\begin{itemize}
    \item \textbf{Bond and capital Eulers \eqref{eq:focB}--\eqref{eq:focK}: a
    $1/g_{t+1}$ appears.} The two $\lambda$'s sit at different dates, so their
    trends do not cancel: $\lambda_{t+1}/\lambda_t =
    (\hat\lambda_{t+1}/\hat\lambda_t)/g_{t+1}$.
    \item \textbf{Labor FOC \eqref{eq:focL}: nothing.} $\lambda_t w_t$ is a
    same-date product of an inverse-trend and a trend variable --- stationary.
    \item \textbf{Education and human-capital FOCs
    \eqref{eq:focE}--\eqref{eq:focH}: nothing.} $H$ is a stationary stock, so its
    shadow value $\lambda^{H}$ is stationary too ($\lambda_{t+1}w_{t+1}L_{t+1}$
    is again an inverse-trend times a trend). This is why the model's
    \texttt{lambda\_H} equation carries no $g$ while the capital Euler does.
    \item \textbf{The HC technology in \eqref{eq:hlom}: the $\Lambda_{t-1}$
    scaling.} Public capital $K^{GE}$ trends while $H$ does not; dividing by the
    trend inside the accumulation technology is what keeps \eqref{eq:hlom}
    balanced. In the model the equation reads \texttt{H =
    (1-deltaH)*H(-1)+chiH*E\^{}gamma*(Kge(-1))\^{}(mu)} because \texttt{Kge} is
    already the detrended stock.
\end{itemize}
See \texttt{../detrending/detrending.tex} for the general logic.

\end{document}
